%\graphicspath{{notes/asy/}}
\thispagestyle{empty}

\title{Math 121B --- Linear Algebra II}
\author{Neil Donaldson}
\date{Winter 2026}
\maketitle 


\def\at#1#2{#1_{#2}}


\boldsubsection{Review of 121A}

We begin by briefly recalling a few basic notions and notations.

\boldinline{Vector Spaces}

Bold-face $\vv$ denotes a \emph{vector} in a \emph{vector space} $V$ over a \emph{field} $\F$. A vector space is closed under \emph{vector addition} and \emph{scalar multiplication}
\[
	\forall \vv_1,\vv_2\in V,\ \forall \lambda_1,\lambda_2\in\F,\quad \lambda_1\vv_1+\lambda_2\vv_2\in V \tag{$\ast$}
\]

\begin{example*}{}
	Make sure you believe ($\ast$) and everything that follows for the vector space 
	\[
		V=\R^2=\{x\vi+y\vj:x,y\in\R\} =\bigl\{\stwovec xy:x,y\in\R\big\}
	\]
	Since the field $\F=\R$ is the real numbers, this is often referred to as a \emph{real} vector space.
\end{example*}


\boldinline{Linear Combinations and Spans} 

Let $\beta\subseteq V$ be a subset of a vector space $V$ over $\F$. A \emph{linear combination} of vectors in $\beta$ is a \emph{finite} sum
  \[
  	\lincom{\lambda}{\vv}{n}
  \]
where $\lambda_j\in\F$ and $\vv_j\in \beta$. The \emph{span} of $\beta$ (over $\F$) is the \emph{subspace} of $V$ comprising all linear combinations:
\[
	\Span_{\F}\beta =\big\{\lincom{\lambda}{\vv}{n}:\lambda_j\in\F,\ \vv_j\in \beta\big\} \le V
\]
  
  
\boldinline{Bases and Co-ordinates} 

A set $\beta\subseteq V$ is a \emph{basis} of $V$ if it has two properties:
\begin{enumerate}
  \item \emph{Linear Independence}. Any linear combination yielding the zero vector is trivial: for distinct $\vv_j\in \beta$,
  \[
  	\lincom{\lambda}{\vv}{n}=\V0\implies \forall j,\lambda_j=0
  \]
  \item \emph{Spanning Set}. $V=\Span_{\F}\beta$: \emph{every} vector in $V$ is a (\emph{finite!}) linear combination of elements of $\beta$,
  \[
  	\forall\vv\in V,\ \exists \lambda_j\in\F \text{ for which }\vv=\lincom\lambda{\vv}n
  \]
\end{enumerate}
  

\begin{thm*}{}
	$\beta$ is a basis of $V\Longleftrightarrow$ \emph{every} $\vv\in V$ is a \emph{unique linear combination} of elements of $\beta$. The coefficients of $\vv$ in this linear combination are its \textbf{co-ordinates} with respect to $\beta$.\smallbreak
	The cardinalities of all basis sets are equal; this value is the \textbf{dimension} $\dim_{\F}V$ of the vector space.
\end{thm*}

It is common to define vector spaces in terms of a basis. For instance, $\R^2$ is defined as the real vector space with basis $\beta=\{\vi,\vj\} =\bigl\{\stwovec 10,\stwovec 01\bigr\}$: this $\beta$ is known as the \emph{standard basis} of $\R^2$.

\begin{examples*}{}{}
	Throughout, let $\F$ be an arbitrary field.\footnotemark{}\par
	\begin{minipage}[t]{0.85\linewidth}\vspace{0pt}
  	\begin{enumerate}
      \item $\F^n=\smash[b]{\Span_{\F}\bigl\{\vect e{n}\big\}}$ is a vector space of dimension $n$, defined in terms of its standard basis: $\ve_k$ is the zero-vector except for a single 1 in the $\smash[t]{k\th}$ entry.\par
      The notation $\vi,\vj,\vk$ is a subscript-avoiding shorthand for $\ve_1,\ve_2,\ve_3$ whenever $n\le 3$.
  	\end{enumerate}
  \end{minipage}
  \hfill
  \begin{minipage}[t]{0.12\linewidth}\vspace{0pt}
  	\hfill$\ve_k=\begin{smatrix}
   		\vdots\\0\\1\\0\\\vdots
    \end{smatrix}$
  \end{minipage} 
  \begin{enumerate}[resume]
%     \item[] Note that $\C^n$ is also a vector space over $\R$: writing $z_k=x_k+iy_k$ where $i^2=-1$, we see that
%     \[
%     	\lincom{z}{\ve}n =\lincom x{\ve}n +\big(\lincom{iy}{\ve}{n}\big)
%     \]
%     whence $\C^n=\Span_{\R}\big\{\vect en,i\ve_1\ldots,i\ve_n\big\}$ and $\dim_{\R}\C^n=2n$. 
    	
    \item $P_n(\F)=\Span_{\F}\bigl\{1,x,x^2,\ldots,x^n\big\}$ is the space of degree $\le n$ \emph{polynomials} with real coefficients. \textbf{Warning}: $\dim_{\F}P_n(\F)=n+1$.\par
    \begin{minipage}[t]{0.77\linewidth}\vspace{-5pt}
    	Suppose $n=2$ and $\F=\R$. Every polynomial is a unique linear combination $p(x)=a+bx+cx^2$ with respect to the standard basis $\beta=\{1,x,x^2\}$. The real numbers $a,b,c$ are the \emph{co-ordinates} of $p$. The \emph{co-ordinate vector} $[p]_\beta$ of $p$ with respect to $\beta$ is the column vector of these co-ordinates.
    \end{minipage}
    \hfill
    \begin{minipage}[t]{0.2\linewidth}\vspace{-5pt}
     	\hfill$[p]_\beta=\threevec{a}{b}{c} \in\R^3$
    \end{minipage}

    \item $P(\F)=\Span_{\F}\bigl\{1,x,x^2,\ldots\big\}$ is the space of polynomials with coefficients in $\F$. This has countably infinite dimension. Note that a polynomial is a \emph{finite sum}: power series are not polynomials! 
    
    \item $C(\R)$ is the real space of continuous functions $f:\R\to\R$. Its dimension is uncountably infinite. Continuity also makes sense over a field in which limits make sense: e.g. $C(\C)$.
  \end{enumerate}
\end{examples*}
  
\footnotetext{%
	 For us, example fields will almost always be $\R$ or $\C$ (sometimes $\Q$). If you've studied abstract algebra, you should also have encountered \emph{extension fields} and \emph{finite fields}. Prior experience with such examples is not necessary for these notes.% No  should be familiar to you, though n abstract algebra, you'll also encounter \emph{finite fields} such as $\F_2=\{0,1\}$ where the operations are taken modulo 2: the vector space $\F_2^2=\bigl\{\stwovec 00,\stwovec 10,\stwovec 01,\stwovec 11\bigr\}$ has 4 elements! 
}


\boldinline{Linearity and Linear Maps} 

A function $\rT:V\to W$ between vector spaces $V$, $W$ over the \emph{same field} $\F$ is \emph{($\F$-)linear} if it respects the linearity properties of both spaces
\[
	\forall \vv_1,\vv_2\in V,\ \forall \lambda_1,\lambda_2\in\F,\quad 	\rT(\lambda_1\vv_1+\lambda_2\vv_2)=\lambda_1\rT(\vv_1)+\lambda_2\rT(\vv_2)
\]
We write $\cL(V,W)$ for the set (vector space!) of linear maps from $V$ to $W$, shortened to $\cL(V)$ when $V=W$. An \emph{isomorphism} is an invertible/bijective linear map.

\begin{thm*}{}
	If $\dim_{\F} V=n$ and $\beta$ is a basis of $V$, then the co-ordinate map $\vv\mapsto[\vv]_\beta$ is an isomorphism of vector spaces $V\to\F^n$.
\end{thm*}


\boldinline{Matrices and Linear Maps} 

Matrix multiplication is a linear map together with a choice of bases.

\begin{example*}{}
	The matrix $A=\begin{smatrix}
	2&-1&-4\\
	0&-1&3
	\end{smatrix}$ defines a linear map $\rL_A:\R^3\to\R^2$ (\emph{left-multiplication by $A$}):
	\[
		\rL_A\threevec xy z =
		\begin{pmatrix}
			2&-1&-4\\
			0&-1&3
		\end{pmatrix}
		\threevec xy z =\twovec{2x-y-4z}{-y+3z}
	\]
	Conversely, the linear map \emph{defines} the matrix $A$: we recover the columns of $A$ by feeding the standard basis vectors $\vi,\vj,\vk$ (of $\R^3$) to the linear map.
	\[
		\twovec 20=\rL_A(\vi)\qquad
		\twovec{-1}{-1}=\rL_A(\vj)\qquad
		\twovec{-4}{3}=\rL_A(\vk)
	\]
\end{example*}

The example applies in general. Suppose $\rT\in\cL(V,W)$ and that $\beta=\{\vv_1,\ldots,\vv_n\}$ and $\gamma=\{\vw_1,\ldots,\vw_m\}$ are bases of $V$, $W$ respectively. Then the \emph{matrix of $\rT$ with respect to $\beta$ and $\gamma$} is the $m\times n$ matrix
\[
	[\rT]_\beta^\gamma
	=\Big([\rT(\vv_1)]_\gamma\ \cdots\ [\rT(\vv_n)]_\gamma\Big)\in M_{m\times n}(\F)
	\tag{$\dag$}
\]
whose $j\th$ column is obtained by feeding the $j\th$ basis vector of $\beta$ to $\rT$ before taking its co-ordinate vector with respect to $\gamma$. This fits naturally with the co-ordinate isomorphisms
\[
	\rT(\vv)=\vw\iff [\rT]_\beta^\gamma[\vv]_\beta=[\vw]_\gamma
\]

\phantomsection\label{pg:changeofcoords}
There are two special cases when $V=W$:
\begin{itemize}
  \item If $\beta=\gamma$, then we simply write $[\rT]_\beta$ instead of $[\rT]_\beta^\beta$.
  \item If $\rT=\I$ is the identity map, then $Q_\beta^\gamma:=[\I]_\beta^\gamma$ is the \emph{change of co-ordinate matrix} from $\beta$ to $\gamma$.
\end{itemize}



\begin{example*}{}{}
	Let $\rT:P_2(\R)\to P_1(\R)$ be the linear map defined by \emph{differentiation}
	\[
		\rT(a+bx+cx^2)=b+2cx
	\]
	The standard bases of $P_2(\R)$ and $P_1(\R)$ are, respectively, $\beta=\{1,x,x^2\}$ and $\gamma=\{1,x\}$. We compute the matrix of $\rT$ with respect to these bases:
	\begin{gather*}
		[\rT(1)]_\gamma=[0]_\gamma=\twovec 00,\qquad 
		[\rT(x)]_\gamma=[1]_\gamma=\twovec 10,\qquad 
		[\rT(x^2)]_\gamma=[2x]_\gamma=\twovec 02\\
		\implies [\rT]_\beta^\gamma
		=\Big([\rT(1)]_\gamma\ [\rT(x)]_\gamma\ [\rT(x^2)]_\gamma\Big)
		=\begin{pmatrix}
			0&1&0\\
			0&0&2
		\end{pmatrix}
	\end{gather*}
	Written in co-ordinates and following $(\dag)$, we observe the original linear map:
	\[
		\left[\rT(a+bx+cx^2)\right]_\gamma 
		\overset{(\dag)}{=} [\rT]_\beta^\gamma [a+bx+cx^2]_\beta 
		=\begin{pmatrix}
			0&1&0\\
			0&0&2
		\end{pmatrix} 
		\threevec ab c 
		=\twovec b{2c} 
		=\left[b+2cx\right]_\gamma
	\]
	Converting between linear maps and matrix multiplication is a central technique of linear algebra. As a quick pre-requisite test, work through the following:\medbreak
	Consider a second basis $\eta=\big\{1+x,x+x^2,x^2+1\big\}$ of $P_2(\R)$.
	\begin{enumerate}
	  \item If $[p]_\eta=\sthreevec 21{-3}$, what is the \emph{polynomial} $p(x)$?
	  \item If $p(x)=2-3x^2$, what is $[p]_\eta$?
	  \item Verify that $\eta$ is indeed a basis of $P_2(\R)$.
		\item Show that $[\rT]_\eta^\gamma =\begin{smatrix}
			1&1&0\\
			0&2&2
		\end{smatrix}$.
		\item Compute the change of co-ordinate matrix $Q_\eta^\beta$ and check that the matrices of $\rT$ are related by
		\[
			[\rT]_\eta^\gamma = [\rT]_\beta^\gamma Q_\eta^\beta
		\]
	\end{enumerate}
\end{example*}


\clearpage


%\setcounter{section}{4}
\section{Diagonalizability \& the Cayley--Hamilton Theorem}

\subsection{Eigenvalues, Eigenvectors \& Diagonalization\ (Review)}

Diagonalization should be mostly familiar, though it is worth a more thorough review.

\begin{defn}{}{}
	Suppose $V$ is a vector space over $\F$ and $\rT\in\cL(V)$. A \emph{non-zero} $\vv\in V$ is an \emph{eigenvector} of $\rT$ with \emph{eigenvalue} $\lambda\in\F$ (together $(\lambda,\vv)$ are an \emph{eigenpair}) if
	\[
		\rT(\vv)=\lambda\vv
	\]
	The \emph{eigenspace} of an eigenvalue $\lambda$ is the nullspace $E_\lambda:=\cN(\rT-\lambda\I)$. Its \emph{geometric multiplicity} is $\dim E_\lambda$.\smallbreak
	We say that $\rT$ is \emph{diagonalizable} if there exists a basis of eigenvectors: an \emph{eigenbasis.}
\end{defn}



\boldsubsubsection{Eigenvalues and Eigenvectors in finite dimensions}

A simple diagonalizability criterion may be proved straightforwardly by induction.

\begin{lemm}{}{evectindep}
	Eigenvectors $\vect{\vv}{k}$ with distinct eigenvalues are linearly independent.\smallbreak
	If $\dim_{\F} V=n$ and $\rT\in\cL(V)$ has $n$ distinct eigenvalues, then $\rT$ is diagonalizable.
\end{lemm}

In finite dimensions, we often work with matrix examples or by choosing bases so as to view a linear maps as matrix multiplication. As such:
\begin{itemize}
	\item The eigenvalues/vectors of a square matrix $A\in M_n(\F)$ are precisely those of the linear map $\rL_A\in\cL(\F^n)$, \emph{left-multiplication} by $A$ (that is, $\rL_A(\vx)=A\vx$).
	\item If $\dim_{\F} V=n$ and $\epsilon$ is a basis, then $\rT(\vv)=\lambda\vv \Longleftrightarrow [\rT]_\epsilon[\vv]_\epsilon=\lambda[\vv]_\epsilon$. To be diagonalizable means there is a basis of eigenvectors $\beta=\{\vect vn\}$ such that $[\rT]_\beta$ is \emph{diagonal}
	\[
		[\rT]_\beta=
		\begin{pmatrix}
			\lambda_1&0&\cdots&0\\
			0&\lambda_2&&\vdots\\
			\vdots&&\ddots&0\\
			0&\cdots&0&\lambda_n
		\end{pmatrix}
	\]
\end{itemize}

In such situations, there is systematic method for computing eigenvalues and eigenvectors:
\begin{enumerate}
  \item (If necessary) \ Choose any basis $\epsilon$ of $V$ and compute the matrix $A=[\rT]_\epsilon\in M_n(\F)$.
  \item Observe that
  \begin{align*}
	  \lambda\in\F\text{ is an eigenvalue }&\iff \exists [\vv]_\epsilon\in\F^n\setminus\{\V0\}\text{ such that }A[\vv]_\epsilon=\lambda[\vv]_\epsilon\\
	  &\iff \exists [\vv]_\epsilon\in\F^n\setminus\{\V0\}\text{ such that }\left(A-\lambda I\right)[\vv]_\epsilon=\V0\\
	  &\iff \det\left(A-\lambda I\right)=0
  \end{align*}
  This last is a degree $n$ polynomial equation whose roots are the eigenvalues.
  \item For each eigenvalue $\lambda_j$, compute the eigenspace $E_{\lambda_j}=\cN(\rT-\lambda_j\I)$ to find the eigenvectors. Remember that $E_{\lambda_j}$ is a subspace of the original vector space $V$, so translate back if necessary!
\end{enumerate}



\begin{defn}{}{}
	The \emph{characteristic polynomial} of $\rT\in\cL(V)$ is the degree-$n$ polynomial\footnotemark{}
	\[
		p(t):=\det(\rT-t\I)
	\]
	The eigenvalues of $\rT$ are precisely the solutions to the \emph{characteristic equation} $p(t)=0$.
\end{defn}

\footnotetext{%
	The determinant may be computed by considering the matrix of $\rT$ with respect to \emph{any} basis: $p(t)$ is independent of this choice (Exercise \ref{exs:detTbasis}).%
}


\begin{examples}{}{intro}
	\exstart $A=
	\begin{smatrix}
		0&-1\\1&0
	\end{smatrix}$ has characteristic polynomial $p(t)=t^2+1=(t+i)(t-i)$. Viewed as a linear map $\rL_A\in\cL(\R^2)$, this matrix has \emph{no eigenvalues} and \emph{no eigenvectors}!
	
	\begin{enumerate}\setcounter{enumi}{1}
	  \item[] However, viewed as a linear map $\rL_A\in\cL(\C^2)$, the matrix has two eigenvalues $\pm i$. Indeed
		\[
			(A-i\I)\vv=
			\begin{pmatrix}
	  		-i&-1\\
	  		1&-i
	  	\end{pmatrix}
	  	\vv\implies E_i=\Span\twovec{i}1
	  \]
		Similarly $E_{-i}=\Span\stwovec{-i}1$. We therefore have an eigenbasis $\beta=\bigl\{\stwovec{i}1,\stwovec{-i}1\bigr\}$ (of $\C^2$), with respect to which the linear map has a diagonal matrix:
		\[
			[\rL_A]_\beta=
			\begin{pmatrix}
				i&0\\
				0&-i
			\end{pmatrix}
		\]
	  \item\label{ex:intro2} Let $\rT\in\cL\big(P_2(\R)\big)$ be defined by
		\[
			\rT(f)(x)=f(x)+(x-1)f'(x)
		\]
		With respect to the standard basis $\epsilon=\{1,x,x^2\}$, the linear map has a non-diagonal matrix:
		\[
			A=[\rT]_\epsilon=
			\begin{pmatrix}
			1&-1&0\\
			0&2&-2\\
			0&0&3
		\end{pmatrix}
		\implies p(t)=\det(A-t I)=(1-t)(2-t)(3-t)
	\]
	With three distinct eigenvalues, Lemma \ref{lemm:evectindep} says that $\rT$ is diagonalizable. To find the eigenvectors, we compute the nullspaces.
	\begin{description}
	  \item[\normalfont$\lambda_1=1$:]\ $\V0=(A-\lambda_1 I)[\vv_1]_\epsilon=
	  \begin{smatrix}
			0&-1&0\\
			0&1&-2\\
			0&0&2
	  \end{smatrix}
	  [\vv_1]_\epsilon
	  \Longrightarrow [\vv_1]_\epsilon\in\Span\sthreevec 100
	  \Longrightarrow E_1=\Span\{1\}$
		\item[\normalfont$\lambda_2=2$:]\ $A-\lambda_2 I=
		\begin{smatrix}
			-1&-1&0\\
			0&0&-2\\
			0&0&1
		\end{smatrix}
		\Longrightarrow [\vv_2]_\epsilon\in\Span\sthreevec 1{-1}0
		\Longrightarrow E_2=\Span\{1-x\}$
		\item[\normalfont$\lambda_3=3$:]\ $A-\lambda_3 I=
		\begin{smatrix}
			-2&-1&0\\
			0&-1&-2\\
			0&0&0
		\end{smatrix}
		\Longrightarrow [\vv_3]_\epsilon\in\Span\sthreevec 1{-2}1
		\Longrightarrow E_3=\Span\{1-2x+x^2\}$
	\end{description}
	A sensible choice of eigenvectors produces an eigenbasis, with respect to which the linear map is necessarily diagonal
	\begin{gather*}
		\beta =\{\vv_1,\vv_2,\vv_3\} =\big\{1,1-x,1-2x+x^2\big\} 
		=\big\{1,1-x,(1-x)^2\big\}\\[5pt]
		\rT\bigl(a+b(1-x)+c(1-x)^2\bigr)=a+2b(1-x)+3c(1-x)^2\\
		[\rT]_\beta=
		\begin{pmatrix}
			1&0&0\\
			0&2&0\\
			0&0&3
		\end{pmatrix}
	\end{gather*}
	\end{enumerate}
\end{examples}


\boldsubsubsection{Conditions for diagonalizability of finite-dimensional operators}

We borrow a little terminology from the theory of polynomials.

\begin{defn}{}{}
	Let $\F$ be a field and $p(t)$ a polynomial with coefficients in $\F$.
	\begin{enumerate}
	  \item Let $\lambda\in\F$ be a root: $p(\lambda)=0$. The \emph{algebraic multiplicity} $\mult(\lambda)$ is the largest power of $\lambda-t$ to divide $p(t)$. Otherwise said, there exists some polynomial $q(t)$ such that
	  \[
	  	p(t)=(\lambda-t)^{\mult(\lambda)}q(t)\quad \text{and}\quad q(\lambda)\neq 0
	  \]
	  \item $p(t)$ \emph{splits} over $\F$ if it factorizes completely into linear factors: $\exists a,\lambda_1,\ldots,\lambda_k\in\F$ such that
	  \[
	  	p(t)=a(\lambda_1-t)^{m_1}\cdots(\lambda_k-t)^{m_k} \tag{$\ast$}
	  \]
	  When $p(t)$ splits, the algebraic multiplicities sum to the degree $n$ of the polynomial
	  \[
	  	n=m_1+\cdots +m_k
	  \]\vspace{-6pt}
	\end{enumerate}
\end{defn}


We are most interested when $p(t)=\det(\rT-t\I)$ is the \emph{characteristic polynomial} of a linear map $\rT\in\cL(V)$. If such a polynomial splits, then (with notation as in $(\ast)$) two things are true:
\begin{itemize}
  \item $a=1$, by considering the degree $n$ term of $p(t)=(-1)^nt^n+\cdots$.
  \item $\lambda_1,\ldots,\lambda_k$ are precisely the distinct eigenvalues of $\rT$.
\end{itemize}

\begin{example*}{\ref{ex:intro}.1, cont.}{}
	The field matters! For instance $p(t)=t^2+1=(t-i)(t+i)=-(i-t)(-i-t)$ splits over $\C$ but not over $\R$. Its roots are plainly $\pm i$.
\end{example*}

For the purposes of review, here is the main result. We will prove it in the next section.

\begin{thm}{}{diagcrit}
	Let $V$ be finite-dimensional. A linear map $\rT\in\cL(V)$ is diagonalizable if and only if two conditions both hold:
	\begin{enumerate}
	  \item Its characteristic polynomial splits over $\F$.
	  \item The geometric and algebraic multiplicities of each eigenvalue are equal; $\dim E_{\lambda_j}=\operatorname{mult}(\lambda_j)$.
	\end{enumerate}
\end{thm}

\begin{example}{}{invariantex}
	The matrix $A=
	\begin{smatrix}
		3&1&0\\
		0&3&0\\
		0&0&5
	\end{smatrix}$
	is easily seen to have eigenvalues $\lambda_1=3$ and $\lambda_2=5$. Indeed
	\begin{gather*}
		p(t)=(3-t)^2(5-t),\qquad \mult(3)=2,\qquad \mult(5)=1\\[5pt]
		E_3=\Span\sthreevec 100,\qquad E_5=\Span\sthreevec 001,\qquad \dim E_3=\dim E_5=1
	\end{gather*}
	This matrix is \emph{non-diagonalizable} since $\dim E_3=1\neq 2=\operatorname{mult}(3)$.
\end{example}

\bigskip

Everything prior to this \emph{should} be review. If it feels very unfamiliar, revisit your notes from 121A, particularly sections 5.1 and 5.2 of the textbook.



\begin{exercises}
	\exstart For each matrix over $\R$: find its characteristic polynomial, its eigenvalues/spaces, and its algebraic and geometric multiplicities. Finally, decide if it is diagonalizable.
	\begin{enumerate}\setcounter{enumi}{1}
  	\item[]\begin{enumerate}
    	\item $A=
    	\begin{smatrix}
	  		2&0&0&0\\
	  		0&3&1&0\\
	  		0&0&3&1\\
	  		0&0&0&3
  		\end{smatrix}$\qquad\quad
  		(b)\ \  $B=
  		\begin{smatrix}
	  		-1&6&0&0\\
	  		-2&6&0&0\\
	  		0&0&3&0\\
	  		0&0&0&3
  		\end{smatrix}$
  	\end{enumerate}
  	
  	
	 	\item Suppose $A$ is a \emph{real} matrix with eigenpair $(\lambda,\vv)$. Show that the complex conjugates $(\cl\lambda,\cl\vv)$ also form an eigenpair.
	  
	  
	  \item Show that the characteristic polynomial of $A=
	  \begin{smatrix}
	  	3&-4\\
	  	4&3
	  \end{smatrix}$
	  does not split over $\R$. Diagonalize $A$ over the complex numbers $\C$.
	  
	  
	  \item Give an example of a $2\times 2$ matrix whose entries are \emph{rational numbers} and whose characteristic polynomial splits over $\R$, but not over $\Q$.
	  
	  
	  \item Diagonalize the matrix $A=
	  \begin{smatrix}
	  	2i&1\\
	  	2&0
	  \end{smatrix}$.
	  
	  
	  \item If a polynomial $p(t)$ splits over a field $\F$, explain why
	  \[
	  	\det \rT=\lambda_1^{\mult(\lambda_1)}\cdots\lambda_k^{\mult(\lambda_k)}
	  \]
	  where $\lambda_1,\ldots,\lambda_k$ are the distinct eigenvalues of $\rT$.
	  
	  
	  \item Suppose $\rT\in\cL(V)$ is invertible with eigenvalue $\lambda$. Prove that $\lambda^{-1}$ is an eigenvalue of $\rT^{-1}$ with the same eigenspace $E_\lambda$. If $\rT$ is diagonalizable, prove that $\rT^{-1}$ is also diagonalizable.\par
	  \emph{Your argument should work even if $V$ is infinite-dimensional: no matrices!}
	  	
	  	
	  \item\label{exs:detTbasis} Throughout this question, suppose $\rT\in\cL(V)$ is a linear map on a finite-dimensional vector space $V$, and that $\beta,\gamma$ are bases of $V$. We use the \emph{change of co-ordinate matrix} $Q_\beta^\gamma$, defined as follows:
	  \[
	  	\forall \vv\in V,\ [\vv]_\gamma=Q_\beta^\gamma[\vv]_\beta
	  \]
	  \begin{enumerate}
	    \item Prove that
	  	\[
	  		[\rT]_\gamma=Q_\beta^\gamma [\rT]_\beta (Q_\beta^\gamma)^{-1}
	  	\]
	  	\item Define $\det\rT:=\det[\rT]_\beta$, where $\beta$ is \emph{any} basis of $V$. Explain why the choice of basis does not matter: that is, if $\gamma$ is any other basis of $V$, then $\det[\rT]_\gamma=\det[\rT]_\beta$.
	  	\item Suppose $A\in M_n(\F)$ is diagonalizable with eigenbasis $\beta=\{\vect vn\}$. Prove that there exists a matrix $M$ such that
	  	\[
	  		A=MDM^{-1}
	  	\]
	  	where $D=\diag(\lambda_1,\ldots,\lambda_n)$ is the diagonal matrix of eigenvalues.
	  \end{enumerate}
	  
	  
	  \item Prove Lemma \ref{lemm:evectindep}.
	\end{enumerate}
\end{exercises}


\clearpage





\subsection{Invariant Subspaces and the Cayley--Hamilton Theorem}

The proof of Theorem \ref{thm:diagcrit} is facilitated by a new concept, of which eigenspaces are a special case.

\begin{defn}{}{}
	Suppose $\rT\in\cL(V)$. A subspace $W$ of $V$ is \emph{$\rT$-invariant} if $\rT(W)\subseteq W$. In such a case, the \emph{restriction} of $\rT$ to $W$ is the linear map
	\[
		\rT_W:W\to W:\vw\mapsto \rT(\vw)
	\]
	As a shorthand, if $A\in M_n(\F)$ then an $\rL_A$-invariant subspace $W\le\F^n$ is simply called \emph{$A$-invariant.}
\end{defn}


\begin{examples}{}{easyinvariantss}
	\exstart For any $\rT\in\cL(V)$, the trivial subspace $\{\V0\}$ and the entirety of $V$ are $\rT$-invariant.
	\begin{enumerate}\setcounter{enumi}{1}
	  \item Every eigenspace of $\rT$ is $\rT$-invariant: if $\vv\in E_\lambda$, then $\rT(\vv)=\lambda\vv\in E_\lambda$.
	  \item\label{ex:easyinv3} (Example \ref{ex:invariantex}, cont.) \ If \smash{$A=
	  \begin{smatrix}
			3&1&0\\
			0&3&0\\
			0&0&5
	  \end{smatrix}
	  \in M_3(\R)$} 
	  then $W=\Span\left\{\vi,\vj\right\}\le\R^3$ is $A$-invariant:
	  \[
	  	A(x\vi+y\vj)=(3x+y)\vi+3y\vj\in W
	  \]
	  Alongside the trivial examples \& eigenspaces, $\Span\{\vi,\vk\}=E_3\oplus E_5$ is also $A$-invariant.
	\end{enumerate}
\end{examples}


To prove the diagonalization criterion (Theorem \ref{thm:diagcrit}), we need to see how to factorize the characteristic polynomial. It turns out that factors of $p(t)$ correspond to $\rT$-invariant subspaces!

\begin{example}{}{}
	$W=\Span\{\vi,\vj\}$ is an invariant subspace of $A=
	\begin{smatrix}
	  1&2&4\\
	  0&3&1\\
	  0&0&2
  \end{smatrix}
  \in M_3(\R)$. With respect to the standard basis, the restriction of $\rL_A$ to $W$ has matrix 
  $\begin{smatrix}
  	1&2\\
  	0&3
  \end{smatrix}$. 
  The characteristic polynomial $p_W(t)$ of the restriction is a factor of the whole,
  \[
  	p(t)=(1-t)(2-t)(3-t)=(2-t)p_W(t)
  \]
\end{example}


\begin{thm}{}{invariantsub}
	Suppose $\rT\in\cL(V)$, that $\dim V$ is finite and that $W$ is a $\rT$-invariant subspace of $V$. Then the characteristic polynomial $p_W(t)$ of the restriction $\rT_W$ divides that of $\rT$.
\end{thm}


The proof simply abstracts the approach of the example.


\begin{proof}
	Extend a basis $\beta_W$ of $W$ to a basis $\beta$ of $V$. Since $\rT(\vw)\in\Span\beta_W$ for each $\vw\in W$, we see that the matrix of $[\rT]$ has block form
  \[
  	[\rT]_\beta=\left(\!
  	\begin{array}{c|c}
  		A&B\\\hline
  		O&C
  	\end{array}
  	\!\right)
  	\implies p(t)=\det(A-tI)\det(C-tI)=p_W(t)\det(C-tI)
  \]
  where $p_W(t)$ is the characteristic polynomial, and $A=[\rT_W]_{\beta_W}$ the matrix of the restriction $\rT_W$.
\end{proof}


\begin{cor}{}{einvariant}
	If $\lambda$ is an eigenvalue of $\rT$, then $\rT_{E_\lambda}=\lambda\I_{E_\lambda}$ is a multiple of the identity. In particular:
	\begin{enumerate}
	  \item The characteristic polynomial of the restriction $\rT_{E_\lambda}$ is $p_\lambda(t)=(\lambda-t)^{\dim E_\lambda}$.
	  \item Since $p_\lambda(t)$ divides the characteristic polynomial of $\rT$, we see that $\dim E_\lambda\le\operatorname{mult}(\lambda)$.
	\end{enumerate}
\end{cor}


We are now in a position to state and prove an extended version of Theorem \ref{thm:diagcrit}.

\begin{thm}{}{}
	Suppose $\dim_{\F} V=n$ and that $\rT\in\cL(V)$ has distinct eigenvalues $\lambda_1,\ldots,\lambda_k$. The following are equivalent:
  \begin{enumerate}
    \item $\rT$ is diagonalizable.
    \item The characteristic polynomial splits over $\F$ and $\dim E_{\lambda_j}=\operatorname{mult}(\lambda_j)$ for each $j$, indeed
    \[
    	p(t)=p_{\lambda_1}(t)\cdots p_{\lambda_k}(t)
    	=(\lambda_1-t)^{\dim E_{\lambda_1}} \cdots (\lambda_k-t)^{\dim E_{\lambda_k}}
    \]
    \item \smash{$\sum\limits_{j=1}^k\dim E_{\lambda_j}=n$}
    \item $V=E_{\lambda_1}\oplus \cdots\oplus E_{\lambda_k}$
	\end{enumerate}
\end{thm}


\begin{example}{}{}
	The data for the diagonalizable matrix 
	\smash[b]{$A=
	\begin{smatrix}
	  7&0&-12\\
	  0&1&0\\
	  2&0&-3
	\end{smatrix}$}
	is shown in the table.\par
	\begin{minipage}[t]{0.56\linewidth}\vspace{0pt}
		\begin{itemize}
		  \item $p(t)=(1-t)^2(3-t)$ splits
		  \item $\R^4=E_1\oplus E_3$.
		  \item $\beta=\left\{\sthreevec 201,\sthreevec 010,\sthreevec 301\right\}$ is an eigenbasis, with respect to which the map is diagonal $[\rL_A]_\beta=
			\begin{smatrix}
			  1&0&0\\
			  0&1&0\\
			  0&0&3
		  \end{smatrix}$
		\end{itemize}
	\end{minipage}
	\hfill
	\begin{minipage}[t]{0.42\linewidth}\vspace{0pt}
		\hfill
		\def\arraystretch{1.45}
		$\begin{array}{c|c|c}
			\lambda&1&3\\\hline
			\operatorname{mult}(\lambda)&2&1\\\hline
			E_\lambda&\Span\left\{\sthreevec 201,\sthreevec 010\right\}&\Span\sthreevec 301\\\hline
			\dim E_\lambda&2&1
		\end{array}$
	\end{minipage}
\end{example}


\begin{proof}
	($1 \Rightarrow 2$) \ If $\rT$ is diagonalizable with eigenbasis $\beta$, then $[\rT]_\beta$ is diagonal. But then\vspace{-22pt}
	\begin{description}
	  \item[] 
	  \[
	  	p(t) =\det\bigl(\diag(\lambda_1-t,\ldots,\lambda_k-t)\bigr) 
	  	=(\lambda_1-t)^{m_1}\cdots (\lambda_k-t)^{m_k}
	  \]
	  splits and $\sum\operatorname{mult}(\lambda_i)=n$. Plainly the cardinality $n$ of an eigenbasis is \emph{at most} $\sum\dim E_{\lambda_i}$. Corollary \ref{cor:einvariant} supplies the second inequality in the following:
	  \[
	  	n\le\sum\dim E_{\lambda_j}\le\sum\operatorname{mult}(\lambda_j)=n
	  \]
	  All inequalities are necessarily \emph{equalities}, whence $\dim E_{\lambda_j}=\operatorname{mult}(\lambda_j)$ for each $j$.
	  
	  \item[\normalfont{($2\Rightarrow 3$)}] $p(t)$ splits $\Longrightarrow n=\sum\operatorname{mult}(\lambda_j)=\sum\dim E_{\lambda_j}$
	  
	  \item[\normalfont{($3\Rightarrow 4$)}] We sketch an induction argument.\smallbreak
	  Given $j\le k$, assume $E_{\lambda_1}\oplus\cdots\oplus E_{\lambda_j}$ exists. If $(\lambda_{j+1},\vv_{j+1})$ is an eigenpair, then $\vv_{j+1}\not\in E_{\lambda_1}\oplus\cdots\oplus E_{\lambda_j}$ for otherwise Lemma \ref{lemm:evectindep} is contradicted. We conclude that $E_{\lambda_1}\oplus\cdots\oplus E_{\lambda_{j+1}}$ is a direct sum.\smallbreak
	  By induction, $E_{\lambda_1}\oplus\cdots\oplus E_{\lambda_k}$ exists. By assumption (condition 3) this direct sum has dimension $n=\dim V$ and therefore \emph{equals} $V$.
	  
	  \item[\normalfont{($4\Rightarrow 1$)}] For each $j$, choose a basis $\beta_j$ of $E_{\lambda_j}$. Then $\beta:=\beta_1\cup\cdots\cup \beta_k$ is easily seen to be a basis of $V$ (this is essentially condition 4, our assumption). Since every element of $\beta$ is an eigenvector of $\rT$, we have an eigenbasis and so $\rT$ is diagonalizable.\qedhere
	\end{description}
\end{proof}




\boldsubsubsection{$\rT$-cyclic Subspaces and the Cayley--Hamilton Theorem}

We introduce a general family of invariant subspaces before employing them to prove a startling result.

\begin{defn}{}{}
	Let $\rT\in\cL(V)$ and let $\vv\in V$. The \emph{$\rT$-cyclic subspace} generated by $\vv$ is the span
	\[
		\ip\vv=\Span\bigl\{\vv,\rT(\vv),\rT^2(\vv),\ldots\big\}
	\]\vspace{-5pt}
\end{defn}

As you should expect by now, we say $A$-cyclic if $\rT=\rL_A$ is left-multiplication by a matrix.


\begin{example*}{\ref{ex:invariantex}, \ref*{ex:easyinvariantss}.\ref{ex:easyinv3}, cont.}{}
	Let $A=
	\begin{smatrix}
		3&1&0\\
		0&3&0\\
		0&0&5
	\end{smatrix}$, and $\vv=\vi+\vk$. It is easy to see that
	\[
		A\vv=3\vi+5\vk,\quad A^2\vv=9\vi+25\vk,
		\quad\ldots,\quad 
		A^m\vv=3^m\vi+5^m\vk
	\]
	Plainly $\Span\{\vi,\vk\}=\ip{\vi+\vk}$ is a $A$-cyclic subspace.
\end{example*}


The proof of the following is left as an exercise.

\begin{lemm}{}{cyclicbasic}
	The $\rT$-cyclic subspace $\ip\vv$ is the smallest $\rT$-invariant subspace of $V$ containing $\vv$. More specifically:
	\begin{enumerate}
	  \item $\ip\vv$ is $\rT$-invariant.
	  \item If $W\le V$ is $\rT$-invariant and $\vv\in W$, then $\ip\vv\le W$.
	  \item $\dim\ip\vv=1\iff\vv$ is an eigenvector of $\rT$.
	\end{enumerate}
\end{lemm}

We were lucky in the example that the general form $A^m\vv$ was so clear. It is helpful to develop a general approach for identifying the dimension and a basis of a $\rT$-cyclic subspace.\medbreak

Suppose $\vv\neq\V0$ and that $\ip{\vv}$ has finite dimension.\footnote{%
	Necessarily the situation if $\dim V<\infty$, whenever we are thinking about characteristic polynomials\ldots%
}
Let $k\ge 1$ be maximal such that the set
\[
	\big\{\vv,\rT(\vv),\ldots,\rT^{k-1}(\vv)\big\}
\]
is linearly independent.
\begin{itemize}
  \item If $k$ doesn't exist, the \emph{infinite} linearly independent set $\big\{\vv,\rT(\vv),\ldots\big\}$ contradicts $\dim\ip{\vv}<\infty$.
  \item By the maximality of $k$, $\rT^k(\vv)\in\Span\bigl\{\vv,\rT(\vv),\ldots,\rT^{k-1}(\vv)\big\}$. By induction this extends to
  \[
  	j\ge k\implies \rT^j(\vv)\in\Span\bigl\{\vv,\rT(\vv),\ldots,\rT^{k-1}(\vv)\big\}
  \]
\end{itemize}

It follows that $\ip{\vv}=\Span\bigl\{\vv,\rT(\vv),\ldots,\rT^{k-1}(\vv)\big\}$, and we've proved a useful criterion.

\begin{thm}{}{finitetcyclic}
	Suppose $\vv\neq \V0$, then
	\begin{align*}
		\dim\ip{\vv}=k&\iff \big\{\vv,\rT(\vv),\ldots,\rT^{k-1}(\vv)\big\}\text{ is a basis of $\ip{\vv}$}\\
		&\iff k\text{ is maximal such that }\big\{\vv,\rT(\vv),\ldots,\rT^{k-1}(\vv)\big\}\text{ is linearly independent}\\
		&\iff k\text{ is minimal such that }\rT^k(\vv)\in\Span\bigl\{\vv,\rT(\vv),\ldots,\rT^{k-1}(\vv)\big\} 
	\end{align*}
\end{thm}


\begin{examples}{}{finitetcyclic}
	\exstart According to the Theorem, in our ongoing Example could simply have noticed:\vspace{-2pt}
	\begin{enumerate}\setcounter{enumi}{1}
	  \item[]\begin{itemize}
	  	\item $\vv=\vi+\vk$ and $A\vv=3\vi+5\vk$ are linearly independent.
	  	\item $A^2(\vi+\vk)=9\vi+25\vk\in\Span\{\vv,A\vv\}$.
		\end{itemize}
		We'd then conclude that $\ip{\vv}=\Span\{\vv,A\vv\}$ has dimension 2.
		
	  \item\label{ex:finitecyclic2} Let $\rT\big(p(x)\big)=3p(x)-p''(x)$ viewed as a linear map $\rT\in\cL\big(P_2(\R)\big)$ and consider the $\rT$-cyclic subspace generated by the polynomial $p(x)=x^2$
		\[
			\rT(x^2)=3x^2-2,\quad \rT^2(x^2)=\rT(3x^2-2)=3(3x^2-2)-6=9x^2-12,\quad\text{etc.}
		\]
		Certainly $\big\{x^2,\rT(x^2)\big\}$ is linearly independent, but
		\[
			\rT^2(x^2)=9x^2-12 
			=-\textcolor{blue}{9}x^2+\textcolor{red}{6}(3x^2-2)
			\in\Span\bigl\{x^2,\rT(x^2)\big\}
		\]
		We conclude that $\dim\ip{x^2}=2$. An alternative basis for $\ip{x^2}$ is plainly $\{1,x^2\}$.
	\end{enumerate}
\end{examples}

The crux of what follows is a simple observation: when restricting a linear map $\rT$ to a $\rT$-cyclic subspace, the coefficients of the characteristic polynomial and the minimal linear combination in the Theorem coincide.

\begin{tcolorbox}[exstyle]
	Continuing the Example, if $W=\ip{x^2}$ and $\beta_W =\big\{x^2,\rT(x^2)\big\} =\big\{x^2,3x^2-2\big\}$, then
	\[
		[\rT_W]_{\beta_W}=
		\begin{pmatrix}
			0&-\textcolor{blue}{9}\\
			1&\textcolor{red}{6}
		\end{pmatrix} 
		\implies p_W(t)=t^2-\textcolor{red}{6}t+\textcolor{blue}{9}
	\]
\end{tcolorbox}


\begin{lemm}{}{invariantsub2}
	Let $\rT\in\cL(V)$ and suppose $W=\ip\vw$ has $\dim W=k$ and basis
	\[
		\beta_W=\big\{\vw,\rT(\vw),\ldots,\rT^{k-1}(\vw)\big\}
	\]
	in accordance with Theorem \ref{thm:finitetcyclic}. Then:
	\begin{enumerate}
	  %\item $\beta_W=\{\vw,\rT(\vw),\ldots,\rT^{k-1}(\vw)\}$ is a basis of $W$.
	  \item If $\rT^k(\vw)+a_{k-1}\rT^{k-1}(\vw)+\cdots+a_0\vw=\V0$, then the characteristic polynomial of $\rT_W$ is
	  \[
	  	p_W(t)=(-1)^k\left(t^k+a_{k-1}t^{k-1}+\cdots+a_1t+a_0\right)
	  \]
	  \item $p_W(\rT_W)=0$ is the \emph{zero map} on $W$.
	\end{enumerate}
\end{lemm}


\begin{proof}
	\exstart This is a tedious exercise: just expand the relevant determinant row by row\ldots
	\begin{enumerate}\setcounter{enumi}{1}
	% 	\item\begin{align*}
	% 	p_W(t)&=\det([\at\rT W]_{\beta_W}-tI_k)
	% 	=
	% 	\begin{vmatrix}
	% 	-t&0&0&\cdots&0&-a_0\\
	% 	1&-t&0&&0&-a_1\\
	% 	0&1&-t&&0&-a_2\\
	% 	\vdots&&&\ddots&&\vdots\\
	% 	0&0&0&&-t&-a_{k-2}\\
	% 	0&0&0&\cdots&1&-a_{k-1}-t
	% 	\end{vmatrix}
	% 	\\
	% 	&=-t\begin{vmatrix}
	% 	-t&0&&0&-a_1\\
	% 	1&-t&&0&-a_2\\
	% 	\vdots&&\ddots&&\vdots\\
	% 	0&0&&-t&-a_{k-2}\\
	% 	0&0&\cdots&1&-a_{k-1}-t
	% 	\end{vmatrix}
	% 	+(-1)^{k}a_0\begin{vmatrix}
	% 	1&-t&0&\cdots&0\\
	% 	0&1&-t&&0\\
	% 	\vdots&&\ddots&\ddots&\vdots\\
	% 	0&0&&\ddots&-t\\
	% 	0&0&\cdots&&1
	% 	\end{vmatrix}
	% 	\end{align*}
	% 	The second matrix has determinant 1, yielding the $(-1)^ka_0$ term. The first is $-t$ multiplied by a determinant of the same type but one dimension lower. An induction finishes things off.
		
	  \item Write $\rS\in\cL(V)$ for the linear map
		\[
			\rS:=p_W(\rT)=(-1)^k\bigl(\rT^k+a_{k-1}\rT^{k-1}+\cdots+a_0\I\bigr)
		\]
		Part 1 says $\rS(\vw)=\V0$. Since $\rS$ is a polynomial in $\rT$, it commutes with all powers of $\rT$:
		\[
			\forall j,\ \rS\big(\rT^j(\vw)\big) =\rT^j\big(\rS(\vw)\big) =\V0
		\]
		Since $\rS$ is zero on the basis $\beta_W$ of $W$, we see that $\at\rS W$ is the zero function.\qedhere
	\end{enumerate}
\end{proof}


With a little sneakiness, we can drop the $W$'s entirely and observe an intimate relation between a linear map and its characteristic polynomial.

\begin{thm}{Cayley--Hamilton}{}
	If $V$ is finite-dimensional, then $\rT\in\cL(V)$ satisfies its characteristic polynomial: $p(\rT)=0$. 
\end{thm}

\begin{proof}
	Let $\vw\in V$ and consider the cyclic subspace $W=\ip\vw$ generated by $\vw$. By 
	Theorem \ref{thm:invariantsub},
	\[
		p(t)=q_{W}(t)p_{W}(t)
	\]
	for some polynomial $q_{W}$. Lemma \ref{lemm:invariantsub2} says that $\big(p_{W}(\rT)\big)(\vw)=\V0$, whence
	\[
		\big(p(\rT)\big)(\vw)=\V0
	\]
	Since we can apply this reasoning to \emph{any} $\vw\in V$, we conclude that $p(\rT)$ is the zero function.
\end{proof}


\begin{examples}{}{cayley-last}
	\exstart $A=
	\begin{smatrix}
		2&1\\
		3&4
	\end{smatrix}$ has $p(t)=t^2-6t+5$ and we confirm:
	\[
		A^2-6A=
		\begin{pmatrix}
	  	7&6\\
	  	18&19
		\end{pmatrix}
		-6
		\begin{pmatrix}
	  	2&1\\
	  	3&4
		\end{pmatrix}
		=-5I
	\]\vspace{-5pt}
	\begin{enumerate}\setcounter{enumi}{1}
	  \item[] It is a strange approach for this matrix, but the characteristic equation can be used to calculate the inverse of $A$:
	\[
		A^2-6A+5I=0\implies A(A-6I)=-5I\implies A^{-1}=\frac 15(6I-A)=\frac 15
		\begin{pmatrix}
	  	4&-1\\
	  	-3&2
		\end{pmatrix}
	\]
		\item We use Cayley--Hamilton to compute $A^4$ when
		\[
			A=
			\begin{pmatrix}
				2&-1&\frac 83\\
				0&1&-6\\
				0&0&2
			\end{pmatrix}
		\]
		The characteristic polynomial is 
		\[
			p(t)=(2-t)^2(1-t)=4-8t+5t^2-t^3
		\]
		whence
		\begin{align*}
			A^4&=AA^3=A(5A^2-8A+4I)\\
			&=5A^3-8A^2+4A=5(5A^2-8A+4I)-8A^2+4A\\
			&=17A^2-36A+20I=17
			\begin{pmatrix}
				4&-3&\frac{50}3\\
				0&1&-18\\
				0&0&4
			\end{pmatrix}
			-36
			\begin{pmatrix}
				2&-1&\frac 83\\
				0&1&-6\\
				0&0&2
			\end{pmatrix}
			+20
			\begin{pmatrix}
				1&0&0\\
				0&1&0\\
				0&0&1
			\end{pmatrix}\\
			&=
			\begin{pmatrix}
				16&-15&\frac{562}3\\
				0&1&-90\\
				0&0&16
			\end{pmatrix}
		\end{align*}
		
		
		\item\label{exs:cayleyham} (Example \ref*{ex:intro}.\ref{ex:intro2}, cont.) \ The linear map $\rT(f(x))=f(x)+(x-1)f'(x)$ has
		\[
			p(t)=(1-t)(2-t)(3-t)=-t^3+6t^2-11t+6
		\]
		By Cayley--Hamilton, $\rT^3=6\rT^2-11\rT+6\I$. You can check this explicitly, after first computing
		\begin{gather*}
			\rT^2\big(f(x)\big)=f(x)+3(x-1)f'(x)+(x-1)^2f''(x),\ \text{etc.}%\\
			%\rT^3(f(x))=f(x)+7(x-1)f'(x)+6(x-1)^2f''(x)
		\end{gather*}
		We can also apply Cayley--Hamilton to simplify higher powers of $\rT$ and even to compute the inverse:
		\begin{align*}
			\I=\frac 16(\rT^3-6\rT^2+11\rT) \implies
			& \rT^{-1}=\frac 16(\rT^2-6\rT+11\I)\\	
			\implies
			&\rT^{-1}\big(f(x)\big) =f(x)-\frac 12(x-1)f'(x)+\frac 16(x-1)^2f''(x)
		\end{align*}
		This should feel very strange: the inverse of $\rT$ involves more derivatives! 
	\end{enumerate}
\end{examples}


\vfil


\begin{exercises}
	\exstart Given $A=
	\begin{smatrix}
	  3&0&0\\
	  0&2&4\\
	  0&0&2
  \end{smatrix}$ find the $A$-cyclic subspace of $\R^3$ generated by $\vk=\ve_3=\sthreevec 001$.
  
	\begin{enumerate}\setcounter{enumi}{1}
	  \item Given $A=
	  \begin{smatrix}
	  	1&2&4\\
	  	0&3&1\\
	  	0&0&2
	  \end{smatrix}$ and $\vv=\sthreevec 01{-1}$, compute $A(\vv)$ and $A^2(\vv)$. Hence describe the $A$-cyclic subspace $\ip{\vv}$ and its dimension.
	  
	  
		\item Given $A=
		\begin{smatrix}
	  	2&0&0&0\\
	  	0&3&1&0\\
	  	0&0&3&1\\
	  	0&0&0&3
	  \end{smatrix}$, 
	  find \emph{two distinct} $A$-invariant subspaces $W\le\R^4$ such that $\dim W=3$.
	
	
	  \item Suppose $W$ and $X$ are $\rT$-invariant subspaces of $V$. Prove that the sum
	  \[
	  	W+X=\{\vw+\vx:\vw\in W,\vx\in X\}
	  \]
	  is also $\rT$-invariant. 
		
		
		\item Prove all three parts of Lemma \ref{lemm:cyclicbasic}.
	
	
		\item Give an example of an infinite-dimensional vector space $V$, a linear map $\rT\in\cL(V)$, and a vector $\vv$ such that $\ip{\vv}=V$. Explain your answer.
		
		
		\item Consider $\beta=\{\sin x,\cos x,2x\sin x,3x\cos x\}$ and $\rT=\diff x\in\cL(\Span\beta)$. Plainly the subspace $W:=\Span\{\sin x,\cos x\}$ is $\rT$-invariant. Compute the matrices $[\rT]_\beta$ and $[\rT_W]_{\beta_W}$ and observe that
		\[
			p(t)=\bigl(p_W(t)\bigr)^2
		\]
	
		
		\item Verify explicitly that $A=
		\begin{smatrix}
			2&3\\
			0&2
		\end{smatrix}$
		satisfies its characteristic polynomial, in accordance with Cayley--Hamilton.
	  	
	  	
	  \item Check the details of Example \ref*{ex:cayley-last}.\ref{exs:cayleyham}. Moreover:
	  \begin{enumerate}
	    \item Evaluate $\rT^4$ as a linear combination of $\I,\rT$ and $\rT^2$.
	    \item Verify the evaluation of $\rT^{-1}\big(f(x)\big)$.
	  \end{enumerate}
	  
	  
	  \pagebreak[3]
	  
	  
	  \item Suppose $a,b$ are constants with $a\neq 0$ and define $\rT\big(f(x)\big)=af(x)+bf'(x)$.
	  \begin{enumerate}
	    \item Find an expression for the inverse $\rT^{-1}\big(f(x)\big)$ if $\rT\in\cL\big(P_1(\R)\big)$
	    \item Find an expression for the inverse $\rT^{-1}\big(f(x)\big)$ if $\rT\in\cL\big(P_2(\R)\big)$
	  \end{enumerate}
	  \emph{Your answers should be written in terms of $f$ and its derivatives.}
	  
	  
	  \item Let $\rT(f)(x)=f'(x)+\frac 1x\int_0^xf(t)\,\dt$ be a linear map $\rT\in\cL\big(P_2(\R)\big)$.
	  \begin{enumerate}
	    %\item Compute the matrix of $L$ with respect to the standard basis $\epsilon=\{1,x,x^2\}$.
	    \item Find the characteristic polynomial of $\rT$ and identify its eigenspaces. Is $\rT$ diagonalizable?
	    \item Find $a,b,c\in\R$ such that $\rT^3=a\rT^2+b\rT+c\I$.
	    \item What are $\dim\cL\big(P_2(\R)\big)$ and $\dim\Span\{\rT^k:k\in\N_0\}$? Explain.
	  \end{enumerate} 
	  
	  
	  \item If $A=
	  \begin{smatrix}
	  	a&b\\
	  	c&d
	  \end{smatrix}$ 
	  has non-zero determinant, use Cayley--Hamilton to obtain the usual expression for $A^{-1}$.
	  
	  
	  \item Recall Examples \ref{ex:invariantex}, \ref*{ex:easyinvariantss}.\ref{ex:easyinv3} and \ref*{ex:finitetcyclic}.1 with $A=
	  \begin{smatrix}
			3&1&0\\
			0&3&0\\
			0&0&5
	  \end{smatrix}$.
	  \begin{enumerate}
	    \item If $\vv=\sthreevec xyz$, show that $\det\big(\vv,A\vv,A^2\vv\big)=-4y^2z$
			\item Hence determine \emph{all} $A$-cyclic subspaces of $\R^3$.
		\end{enumerate}
		
		
		\item\begin{enumerate}
		  \item Consider Example \ref*{ex:finitetcyclic}.\ref{ex:finitecyclic2} where $\rT\in\cL(P_2(\R))$ is defined by $\rT(p(x))=3p(x)-p''(x)$. Prove that all $\rT$-cyclic subspaces have dimension $\le 2$.
		  \item What if we instead consider $\rS\in\cL(P_2(\R)$ defined by $\rS(p(x))=3p(x)-p'(x)$?
		\end{enumerate}
	  
	  
	  \item We prove part 1 of Lemma \ref{lemm:invariantsub2}.
	  \begin{enumerate}
	    \item Explain why the matrix of $\rT_W$ with respect to the basis $\beta_W$ is
	    \[
	    	[\rT_W]_{\beta_W}=
	    	\begin{pmatrix}
			 		0&0&0&\cdots&0&-a_0\\
			 		1&0&0&&0&-a_1\\
			 		0&1&0&&0&-a_2\\
			 		\vdots&&&\ddots&&\vdots\\
					0&0&0&&0&-a_{k-2}\\
			 		0&0&0&\cdots&1&-a_{k-1}
			 	\end{pmatrix}
			 	\in M_k(\F)
			 \]
	 		\item Compute the characteristic polynomial $p_W(t)=\det\bigl([\at\rT W]_{\beta_W}-tI_k\bigr)$ by expanding the determinant along the first row.
		\end{enumerate}
	\end{enumerate}
\end{exercises}

\iffalse
\fi